\documentclass[12pt]{article}
\usepackage[utf8]{inputenc}
\usepackage[T1]{fontenc}
\usepackage{hyperref}
\title{Paungnyuan Naga Language Resources for K-12 Curriculum Development}
\author{Compiled by Claude Sonnet 5 (claude-sonnet-5) for Andr\'as Kornai}
\date{}
\begin{document}
\maketitle
\section{Introduction}
Paungnyuan Naga (also Makyam, Makyan, Macham Naga, Pongnyun; ISO 639-3 umn) is a Sino-Tibetan language of Sagaing Region, Myanmar. It is spoken in the Lahe and Hkamti townships near the Indian border. Population figures differ: about 4,000 speakers (2008), 4,361 people (2012 survey) and about 4,800 people (2025). Young people increasingly prefer Burmese, and many older speakers also know Lainong.
\section{Primary Educational Resources}
I found no primers, textbooks or school programmes for Paungnyuan.
\section{Linguistic and Reference Materials}
Overview of the language and its classification: \url{https://en.wikipedia.org/wiki/Makyam_language}. The language profile on Joshua Project is at \url{https://joshuaproject.net/index.php/languages/umn}.
\section{Religious and Community Resources}
Mission-oriented people-group profiles are at \url{https://joshuaproject.net/people_groups/22666/BM} and \url{https://www.operation.asiaharvest.org/MYANMAR-Paungnyuan-Naga.pdf}. They give demographic information only. I did not see any scripture in the language in the results.
\section{Digital and Audio Resources}
I found no audio, video or corpus resources.
\section{Curriculum Development Context}
Documentation is thin. The language appears to be stable but is under pressure from Burmese. A curriculum would need community-led orthography and materials development.
\end{document}